% GENERATED FILE. Edit canonical lessons, then run npm run generate:books.
% canonical-source-hash: fnv1a64:6ee971c24f474a46
% canonical-lessons: SA-C508-istaka, SA-C508-cinoti, SA-C508-nirmanam, SA-C508-sandadhati, SA-C508-sudha

\chapter{Bricks and Mortar}
\label{ch:sa-bricks-and-mortar}
\hlchaptermodality{508}

\begin{chapteropening}
\textbf{By the end of this chapter:} \emph{I can talk about building, mending and plastering walls.}

Complete the last lesson of chapter 508: I can talk about building, mending and plastering walls.
\end{chapteropening}

\section[\texorpdfstring{\sk{इष्टका} (iṣṭakā)}{इष्टका (iṣṭakā)}]{\sk{इष्टका} (iṣṭakā) — a brick}

\label{lesson:SA-C508-istaka}



\begin{quote}
Before the new one: say the Sanskrit for blows on a fire, then the Sanskrit for an ornament.
\end{quote}

\subsection*{You'll want to know: \sk{इष्टका}}
\textbf{\sk{इष्टका}} — \emph{iṣṭakā} — \textquotedblleft{}a brick\textquotedblright{}.

\begin{quote}
\sk{कर्मकरः} \sk{इष्टकाः} \sk{आनयति।}
\end{quote}

\begin{quote}
\emph{karmakaraḥ iṣṭakāḥ ānayati.}
\end{quote}

\begin{quote}
The worker brings the bricks.
\end{quote}

A feminine noun.

\begin{cousinweb}
\textbf{\sk{इष्टका}} (\emph{iṣṭakā}) is the word MW records for a brick, first a brick of the sacrificial altar; MW gives no derivation, so its origin is uncertain (MW).
\end{cousinweb}

\begin{grammarlens}[title={What you've built}]
One more word for describing what you lived and felt.
\end{grammarlens}

\subsection*{Your turn}
\begin{itemize}
\raggedright
  \item \emph{Say it:} \emph{iṣṭakā}
  \item \emph{Say it:} \emph{iṣṭakā}, once more
  \item \emph{Say it:} \emph{karmakaraḥ iṣṭakāḥ ānayati.} — The worker brings the bricks.
  \item \emph{Recall:} say the Sanskrit for blows on a fire, then the Sanskrit for an ornament, then say \emph{iṣṭakā} again
\end{itemize}

\begin{tcolorbox}[breakable,colback=teal!4,colframe=teal!35!black,title={Before you move on}]
What does \sk{इष्टका} mean? (\textquotedblleft{}A brick\textquotedblright{}.) Say it once more.
\end{tcolorbox}
\section[\texorpdfstring{\sk{चिनोति} (cinoti)}{चिनोति (cinoti)}]{\sk{चिनोति} (cinoti) — piles up; builds up}

\label{lesson:SA-C508-cinoti}



\begin{quote}
Before the new one: say the Sanskrit for an ornament, then the Sanskrit for a brick.
\end{quote}

\subsection*{You'll want to know: \sk{चिनोति}}
\textbf{\sk{चिनोति}} — \emph{cinoti} — \textquotedblleft{}piles up; builds up\textquotedblright{}.

\begin{quote}
\sk{कर्मकरः} \sk{इष्टकाः} \sk{चिनोति।}
\end{quote}

\begin{quote}
\emph{karmakaraḥ iṣṭakāḥ cinoti.}
\end{quote}

\begin{quote}
The worker stacks the bricks.
\end{quote}

Past (he …): \sk{सः} \sk{भित्तिं} \sk{चितवान्}, \emph{saḥ bhittiṃ citavān}.

\begin{cousinweb}
\textbf{\sk{चिनोति}} (\emph{cinoti}) comes from the root √ci, \textquotedblleft{}to arrange in a row, to pile up, to build\textquotedblright{}; its past participle is \emph{cita} (MW).
\end{cousinweb}

\begin{grammarlens}[title={What you've built}]
One more word for describing what you lived and felt.
\end{grammarlens}

\subsection*{Your turn}
\begin{itemize}
\raggedright
  \item \emph{Say it:} \emph{cinoti}
  \item \emph{Say it:} \emph{cinoti}, once more
  \item \emph{Say it:} \emph{karmakaraḥ iṣṭakāḥ cinoti.} — The worker stacks the bricks.
  \item \emph{Recall:} say the Sanskrit for an ornament, then the Sanskrit for a brick, then say \emph{cinoti} again
\end{itemize}

\begin{tcolorbox}[breakable,colback=teal!4,colframe=teal!35!black,title={Before you move on}]
What does \sk{चिनोति} mean? (\textquotedblleft{}Piles up; builds up\textquotedblright{}.) Say it once more.
\end{tcolorbox}
\section[\texorpdfstring{\sk{निर्माणम्} (nirmāṇam)}{निर्माणम् (nirmāṇam)}]{\sk{निर्माणम्} (nirmāṇam) — construction}

\label{lesson:SA-C508-nirmanam}



\begin{quote}
Before the new one: say the Sanskrit for a brick, then the Sanskrit for piles up.
\end{quote}

\subsection*{You'll want to know: \sk{निर्माणम्}}
\textbf{\sk{निर्माणम्}} — \emph{nirmāṇam} — \textquotedblleft{}construction\textquotedblright{}.

\begin{quote}
\sk{नगरे} \sk{नवस्य} \sk{सेतोः} \sk{निर्माणं} \sk{भवति।}
\end{quote}

\begin{quote}
\emph{nagare navasya setoḥ nirmāṇaṃ bhavati.}
\end{quote}

\begin{quote}
A new bridge is being built in the city.
\end{quote}

A neuter noun.

\begin{cousinweb}
\textbf{\sk{निर्माणम्}} (\emph{nirmāṇam}) is built on \emph{nis}, \textquotedblleft{}out\textquotedblright{}, and the root √mā, \textquotedblleft{}to measure, to form\textquotedblright{}; MW gives \emph{nirmāṇa} as forming, making, building (MW).
\end{cousinweb}

\begin{grammarlens}[title={What you've built}]
One more word for describing what you lived and felt.
\end{grammarlens}

\subsection*{Your turn}
\begin{itemize}
\raggedright
  \item \emph{Say it:} \emph{nirmāṇam}
  \item \emph{Say it:} \emph{nirmāṇam}, once more
  \item \emph{Say it:} \emph{nagare navasya setoḥ nirmāṇaṃ bhavati.} — A new bridge is being built in the city.
  \item \emph{Recall:} say the Sanskrit for a brick, then the Sanskrit for piles up, then say \emph{nirmāṇam} again
\end{itemize}

\begin{tcolorbox}[breakable,colback=teal!4,colframe=teal!35!black,title={Before you move on}]
What does \sk{निर्माणम्} mean? (\textquotedblleft{}Construction\textquotedblright{}.) Say it once more.
\end{tcolorbox}
\section[\texorpdfstring{\sk{सन्दधाति} (sandadhāti)}{सन्दधाति (sandadhāti)}]{\sk{सन्दधाति} (sandadhāti) — mends; puts together}

\label{lesson:SA-C508-sandadhati}



\begin{quote}
Before the new one: say the Sanskrit for piles up, then the Sanskrit for construction.
\end{quote}

\subsection*{You'll want to know: \sk{सन्दधाति}}
\textbf{\sk{सन्दधाति}} — \emph{sandadhāti} — \textquotedblleft{}mends; puts together\textquotedblright{}.

\begin{quote}
\sk{सः} \sk{जीर्णं} \sk{कटं} \sk{सन्दधाति।}
\end{quote}

\begin{quote}
\emph{saḥ jīrṇaṃ kaṭaṃ sandadhāti.}
\end{quote}

\begin{quote}
He mends the old mat.
\end{quote}

Past (he …): \sk{सः} \sk{जीर्णं} \sk{कटं} \sk{संहितवान्}, \emph{saḥ jīrṇaṃ kaṭaṃ saṃhitavān}.

\begin{cousinweb}
\textbf{\sk{सन्दधाति}} (\emph{sandadhāti}) is built on \emph{sam}, \textquotedblleft{}together\textquotedblright{}, and the root √dhā, \textquotedblleft{}to put\textquotedblright{}; MW gives \emph{saṃdhā} as to put together, to join, and also to mend, to restore; its past participle is \emph{saṃhita} (MW).
\end{cousinweb}

\begin{grammarlens}[title={What you've built}]
One more word for describing what you lived and felt.
\end{grammarlens}

\subsection*{Your turn}
\begin{itemize}
\raggedright
  \item \emph{Say it:} \emph{sandadhāti}
  \item \emph{Say it:} \emph{sandadhāti}, once more
  \item \emph{Say it:} \emph{saḥ jīrṇaṃ kaṭaṃ sandadhāti.} — He mends the old mat.
  \item \emph{Recall:} say the Sanskrit for piles up, then the Sanskrit for construction, then say \emph{sandadhāti} again
\end{itemize}

\begin{tcolorbox}[breakable,colback=teal!4,colframe=teal!35!black,title={Before you move on}]
What does \sk{सन्दधाति} mean? (\textquotedblleft{}Mends; puts together\textquotedblright{}.) Say it once more.
\end{tcolorbox}
\section[\texorpdfstring{\sk{सुधा} (sudhā)}{सुधा (sudhā)}]{\sk{सुधा} (sudhā) — lime plaster}

\label{lesson:SA-C508-sudha}



\begin{quote}
Before the new one: say the Sanskrit for construction, then the Sanskrit for mends.
\end{quote}

\subsection*{You'll want to know: \sk{सुधा}}
\textbf{\sk{सुधा}} — \emph{sudhā} — \textquotedblleft{}lime plaster\textquotedblright{}.

\begin{quote}
\sk{भित्तिः} \sk{सुधया} \sk{श्वेता} \sk{भवति।}
\end{quote}

\begin{quote}
\emph{bhittiḥ sudhayā śvetā bhavati.}
\end{quote}

\begin{quote}
The wall turns white with lime plaster.
\end{quote}

A feminine noun.

\begin{cousinweb}
\textbf{\sk{सुधा}} (\emph{sudhā}) is built on \emph{su}, \textquotedblleft{}good\textquotedblright{}, and the root √dhe, \textquotedblleft{}to drink\textquotedblright{}: MW gives \emph{sudhā} first as nectar, the drink of the gods, and also as whitewash, plaster and mortar (MW).
\end{cousinweb}

\begin{grammarlens}[title={What you've built}]
One more word for describing what you lived and felt.
\end{grammarlens}

\subsection*{Your turn}
\begin{itemize}
\raggedright
  \item \emph{Say it:} \emph{sudhā}
  \item \emph{Say it:} \emph{sudhā}, once more
  \item \emph{Say it:} \emph{bhittiḥ sudhayā śvetā bhavati.} — The wall turns white with lime plaster.
  \item \emph{Recall:} say the Sanskrit for construction, then the Sanskrit for mends, then say \emph{sudhā} again
\end{itemize}

\begin{tcolorbox}[breakable,colback=teal!4,colframe=teal!35!black,title={Before you move on}]
What does \sk{सुधा} mean? (\textquotedblleft{}Lime plaster\textquotedblright{}.) Say it once more.
\end{tcolorbox}
